\section{Empirical Strategy}
\label{sec:empirical_strategy}

This section describes the econometric approach used to test the predictions developed in Section~\ref{sec:theory}. We describe the main specification, discuss identification challenges, and outline the four specification variants employed.

\subsection{Main Specification}

We estimate the relationship between frequent losers and tender outcomes using the following baseline specification:

\begin{equation}
  y_{igt} = \beta \, \textit{losers}_{igt} + \alpha_g + \lambda_t + \mathbf{x}'\delta + \varepsilon_{igt},
  \label{eq:main}
\end{equation}

\noindent where $y_{igt}$ is the outcome (log negotiated price, log number of firms, or log number of bids) for purchase $i$ of item $g$ in year $t$; $\textit{losers}_{igt}$ is a binary indicator equal to one if at least one frequent loser participated in the tender; $\alpha_g$ represents item fixed effects that absorb time-invariant item characteristics (baseline price level, product complexity, typical market structure); $\lambda_t$ captures year fixed effects for common macroeconomic shocks; and $\mathbf{x}$ includes controls for procedure type (convite vs.\ preg\~{a}o) and, in some specifications, public buyer unit (PBU) fixed effects.

The coefficient of interest is $\beta$, which measures the average difference in outcomes between tenders with and without frequent losers, conditional on comparing the same item across years and controlling for procedure type and buyer identity.

\subsection{Specification Variants}

We estimate four specifications of increasing stringency:

\begin{enumerate}
  \item \textbf{General:} Item and year fixed effects with a convite dummy.
    \[ y_{igt} = \beta \, \textit{losers}_{igt} + \gamma \, \textit{convite}_{igt} + \alpha_g + \lambda_t + \varepsilon_{igt} \]

  \item \textbf{General + PBU:} Adds PBU fixed effects to control for systematic differences across buying agencies.
    \[ y_{igt} = \beta \, \textit{losers}_{igt} + \gamma \, \textit{convite}_{igt} + \alpha_g + \lambda_t + \delta_p + \varepsilon_{igt} \]

  \item \textbf{Preg\~{a}o only:} Restricts the sample to preg\~{a}o (electronic auction) tenders, with item, year, and PBU fixed effects. The convite dummy is excluded.

  \item \textbf{Convite only:} Restricts the sample to convite (sealed-bid) tenders, with item, year, and PBU fixed effects.
\end{enumerate}

Standard errors are clustered at the item level in all specifications to account for serial correlation within item markets.

\subsection{Identification}

The ideal experiment would randomly assign frequent losers to tenders and compare outcomes. Since frequent loser participation is endogenous, several identification concerns arise:

\paragraph{Selection bias.} Items in markets where cartels operate may differ from items in competitive markets. Item fixed effects address this concern by comparing the \emph{same} item with and without frequent loser participation across different tenders.

\paragraph{Time-varying confounders.} Economic conditions may simultaneously affect cartel activity and tender outcomes. Year fixed effects absorb common aggregate shocks.

\paragraph{Buyer heterogeneity.} Some PBUs may be more susceptible to collusion due to lower monitoring capacity. The General + PBU specification (variant 2) controls for systematic buyer-level differences.

\paragraph{Reverse causality.} One might worry that high-price tenders attract frequent losers, rather than the reverse. We address this concern through a staggered difference-in-differences design (Section~\ref{sec:robustness}) that exploits the timing of frequent losers' \emph{first} entry into an item market. Pre-trends in the event study provide evidence on whether outcomes changed before frequent losers appeared.

\subsection{Joint Interpretation of Results}

A critical feature of our empirical strategy is the \emph{joint interpretation} of results across dependent variables. Under competitive conditions, more participating firms and more bids should reduce the winning price---the standard competitive effect. Observing simultaneously:
\begin{itemize}
  \item higher prices (positive $\beta$ for log price), and
  \item more participants and bids (positive $\beta$ for log firms and log bids)
\end{itemize}
is inconsistent with competition or tacit collusion but consistent with an explicit cartel that inflates participation to avoid detection (Predictions~\ref{pred:prices}--\ref{pred:nbids}).

This joint test provides the key identification argument: no single-equation result is sufficient, but the \emph{combination} of higher prices with inflated competition indicators is a signature of explicit collusion with detection avoidance.

\subsection{Sample}

The estimation sample consists of all item-level transactions in the BEC platform from 2009 to 2019, restricted to convite and preg\~{a}o procedures, for item types in which at least one tender involved a frequent loser. This restriction ensures that identification comes from within-item variation. For price regressions, we further restrict to completed tenders with valid negotiated prices. The resulting sample contains approximately 1.67 million observations.
